\section{Evaluation \& Results}

\subsection{Evaluation Metrics}

\subsection{Results}

\paragraph{Baseline Performance.}
We first evaluate classical models on raw sensor data. KNN achieves the highest performance
(Accuracy 0.6505, Macro-F1 0.6374), while Logistic Regression and SVM perform worse due to limited
feature representation. This highlights the importance of feature engineering for linear models.

\paragraph{Effect of Dropout.}
Dropout corruption consistently degrades performance as severity increases. The degradation is
more pronounced for GYRO than ACCL, suggesting that gyroscope signals are more sensitive to
missing data.

\paragraph{Effect of Drift.}
Drift has a strong negative impact, particularly on ACCL. Linear models degrade significantly
under drift due to sensitivity to distribution shifts, while KNN is comparatively more stable.

\paragraph{Effect of Gain.}
Gain exhibits non-monotonic behavior. When gain $< 1$, performance decreases moderately due to
signal attenuation. However, when gain $> 1$, KNN performance improves substantially, indicating
that amplification can enhance class separability. This suggests that gain does not purely act as
corruption, but can implicitly reweight informative features.

\paragraph{Effect of Bias.}
Bias introduces moderate degradation across all models, with ACCL being more affected than GYRO.
This indicates that constant offsets distort magnitude-based features.
\paragraph{Effect of Stochastic Noise.}
Stochastic noise has minimal impact on performance and occasionally leads to slight improvements.
This suggests that models are inherently robust to small perturbations and may benefit from noise
as a form of regularization.

\paragraph{Effect of Resolution.}
Resolution changes have negligible impact across all models, indicating that activity recognition
relies on coarse motion patterns rather than fine-grained signal precision.

\paragraph{Overall Trends.}
Across all corruption types, we observe that not all perturbations degrade performance
monotonically. While dropout and drift significantly harm performance, stochastic noise and
resolution have minimal impact. Interestingly, gain can even improve performance under certain
conditions. These findings highlight that different corruption types affect models in fundamentally
different ways, and robustness depends both on the corruption type and the model architecture.
